% Copyright 2026 icekylinx and GPT-6 Astra Pro. Apache-2.0.
% Focused curation of tcs_research_checkpoint.zip: manuscripts/partial_transpose.tex.
% Retains the circuit, orientation, local projection, border and factor-entry proofs.
% Existing integral primitive remains a cited canonical repository dependency.
\documentclass[11pt]{article}
\usepackage[margin=1in]{geometry}
\usepackage{amsmath,amssymb,amsthm,mathtools}
\usepackage{booktabs,enumitem,microtype}
\usepackage[colorlinks=true,linkcolor=blue,citecolor=blue,urlcolor=blue]{hyperref}
\usepackage{url}
\setlist{itemsep=2pt,topsep=4pt}
\allowdisplaybreaks
\setlength{\emergencystretch}{2em}

\newtheorem{theorem}{Theorem}[section]
\newtheorem{lemma}[theorem]{Lemma}
\newtheorem{proposition}[theorem]{Proposition}
\newtheorem{corollary}[theorem]{Corollary}
\theoremstyle{definition}
\newtheorem{definition}[theorem]{Definition}
\theoremstyle{remark}
\newtheorem{remark}[theorem]{Remark}

\newcommand{\F}{\mathbb F}
\newcommand{\Z}{\mathbb Z}
\newcommand{\Q}{\mathbb Q}
\newcommand{\PT}{\operatorname{PT\!\text{-}rank}}
\newcommand{\rk}{\operatorname{rank}}
\newcommand{\I}{\operatorname{Id}}
\newcommand{\Bnd}{\mathcal B}
\newcommand{\Ori}{\mathcal O}
\newcommand{\cross}{\mathsf P}
\newcommand{\trsp}{\mathsf T}
\newcommand{\polylog}{\operatorname{polylog}}


\title{Partial-Transpose Decompositions of Ordinary Arithmetic Circuits}
\author{icekylinx and GPT-6 Astra Pro}
\date{November 2, 2026}
\begin{document}
\maketitle
\begin{abstract}
A division-free arithmetic circuit of size $s$ in $d$ blocks of
$n$-by-$n$ variables has paired squarefree coefficient matrix of
partial-transpose rank at most $2^{O(d)}(2nsd)^{O(\sqrt d)}$
over every field. Over every commutative unital ring, the same bound holds
for supplied partial transposes of outer products. The proof combines
the existing normalization-one integral crossing primitive with shared
circuit contexts, canonical orientations and lazy squarefree projection.
The geometric border version and uniform factor-entry construction are
retained. This is a structural upper bound and a compact-factor
construction, not an algorithm for arbitrary coefficient evaluation.
\end{abstract}
\paragraph{Scope and dependency.}
Assume $n\ge2$, $d\ge1$ and $s\ge2$; constants may be arbitrary
specified elements of the coefficient ring. Size counts all binary
arithmetic gates and inputs. The sole nonclassical imported algebraic
primitive is Proposition~\ref{prop:primitive}, proved in the existing
repository's finite and integral crossing manuscripts. No rectangular
matrix multiplication theorem, SAT result, or unproved counting
interface is used. The retained set-multilinear theorem is an intermediate
lemma, rather than a separate catalogue result. The reference to
Rossman--Zhu's shifted tensor only distinguishes encodings.
\section{Partial transpose rank and the integral input}

\subsection{Definitions and factorization conventions}

\begin{definition}
Let $M$ be indexed by $(\mathbf i,\mathbf j)\in[n]^d\times[n]^d$.
For $\kappa\subseteq[d]$, the partial transpose
$M^{\Gamma_\kappa}$ exchanges $i_t$ and $j_t$ for every
$t\in\kappa$. A matrix is \emph{PT-basic} if some partial transpose
has ordinary rank one. Its partial transpose rank is the least
number of PT-basic matrices whose sum is $M$.
\end{definition}

For a field $\F$, write $\overline{\PT}_{\F}(M)$ for its
\emph{geometric border PT-rank}: the least $r$ for which $M$ lies
in the Zariski closure, over an algebraic closure of $\F$, of the
set of matrices of PT-rank at most $r$.  The ambient paired
coordinate spaces are held fixed.  This geometric convention
avoids interpreting Zariski closure only on the finite set of
rational points when $\F$ is finite.  Border statements below
concern this measure, not ordinary PT-rank.

The permitted coordinate pairs are part of the definition. If
several pairs are grouped into one larger pair, only whole-block
transposes are available. Every such transpose is still available
in the finer pairing, so refining a grouping cannot increase
PT-rank.

Over a commutative unital ring, we use a supplied factorization
$A=UV^{\trsp}$ through a free module of a stated dimension. This
avoids field-dependent uses of rank. Over fields, such a
factorization immediately implies the usual rank bound. All
ring-uniform statements below retain the factorizations themselves.

For $a,b\ge1$, write $\cross_{\F}(a,b)$ for the PT-rank of
$I_a\boxtimes I_b$ with its two indicated coordinate blocks.
Equivalently, it is the minimum number of rank-one terms needed to
write the crossing
\[
 \delta_{x,v}\delta_{y,u},\qquad
 x,v\in[a],\quad y,u\in[b],
\]
using the two flattenings
\[
 (x,y)\mid(v,u)
 \qquad\text{and}\qquad
 (x,u)\mid(v,y).
\]
The exchange of $u,v$ within a displayed column index is only a
permutation of the columns and does not affect the rank.

\subsection{The supplied normalization-one crossing theorem}

Let $p(j)$ be the partition function, with $p(j)=0$ for $j<0$.
We use the following exact result from the user's repository
\cite{integralcrossing}.

\begin{proposition}[Integral crossing primitive]
\label{prop:primitive}
For every $h,E\ge1$, there are integer tensors
$A_s(x,y,u,v)$, indexed by $-h-E\le s\le h+E$, such that $x,v$
range over the partitions of $h$, $y,u$ range over the partitions
of $E$, and
\begin{equation}
\label{eq:primitive}
 \sum_{s=-h-E}^{h+E}A_s(x,y,u,v)
 =\delta_{x,v}\delta_{y,u}.
\end{equation}
Both flattenings of each piece have explicit integer
factorizations. Their dimensions are at most
\begin{align}
 R_s^{\rm row}&=\sum_{i=0}^{2h}p(E+s-h+i),\label{eq:rowrank}\\
 R_s^{\rm col}&=\sum_{i=0}^{2h}p(E+h-s-i).
 \label{eq:colrank}
\end{align}
The normalizing scalar in \eqref{eq:primitive} is one. The identity
and the supplied factorizations therefore specialize to every
commutative unital ring.
\end{proposition}

The source constructs these pieces by using the integral
symmetric-function lattice, the Schur basis, and its unimodular
Hall pairing in a finite coefficient formulation of the Heisenberg
mode identities. The source also gives an explicit algorithm for
the coefficients. We use the finite identity and the factorization
dimensions, rather than a SAT algorithm or an unconstructed
coefficient oracle.

The pinned source is the following commit of
\texttt{icekylinx/algorithm-frontier}:
\begin{center}
\texttt{2658b60c9d032611492792718bd9d39f721de09f}.
\end{center}
The theorem appears in
\nolinkurl{problems/noncommutative-crossing/manuscripts/integral-pieces.tex}.
The neighboring files \texttt{finite-crossing.tex} and
\texttt{coefficient-computation.tex} supply the finite operator
identities and coefficient computation. The underlying crossing
construction is credited there to the Anthropic manuscript
\cite{anthropic}; the normalization-one integral extension is
credited to the repository authors.

\subsection{A finite profile with a logarithmic estimate}

For $a\ge1$, define
\[
 h(a)=\min\{h\ge1:p(h)\ge a\}.
\]
For $H\ge0$, put
\begin{equation}
\label{eq:Bdef}
 \Bnd(H)=\sum_{j=0}^{H}\bigl(2(H-j)+1\bigr)p(j).
\end{equation}

\begin{lemma}[Finite two-block profile]
\label{lem:finite-profile}
For all $a,b\ge1$, the identity on $[a]\times[b]$ has an integer
decomposition
\begin{equation}
\label{eq:two-group}
 I_a\boxtimes I_b=C_0+C_1^{\Gamma_B},
\end{equation}
where $C_0,C_1$ have explicit integer factorizations whose
dimensions sum to at most
\[
 W(a,b):=\Bnd(h(a)+h(b)).
\]
In particular, each dimension is at most $W(a,b)$, and
$\cross_{\F}(a,b)\le W(a,b)$ over every field.
\end{lemma}

\begin{proof}
Use Proposition~\ref{prop:primitive} with $h=h(a)$ and $E=h(b)$,
and restrict the two label sets to $a$ and $b$ labels. Write
$H=h+E$. For $s\le0$, use the row factorization; for $s>0$, use
the column factorization. The selected dimension is at most
\[
 \sum_{j=0}^{H-|s|}p(j).
\]
Expanding these factorizations into coordinate products and
grouping them by their two orientations gives
\eqref{eq:two-group}. Their total dimension is at most
\[
 \sum_{s=-H}^{H}\sum_{j=0}^{H-|s|}p(j)
 =\sum_{j=0}^{H}\bigl(2(H-j)+1\bigr)p(j)=\Bnd(H).
\]
All factors are integer factors before specialization.
\end{proof}

\begin{lemma}[Summing the partition profile]
\label{lem:B-asymptotic}
There is an absolute constant $C$ such that, for all $H\ge2$,
\begin{equation}
\label{eq:B-upper}
 \Bnd(H)\le CHp(H)
 \le C\exp\!\left(\pi\sqrt{2H/3}\right).
\end{equation}
\end{lemma}

\begin{proof}
The Hardy--Ramanujan asymptotic \cite{dlmf} states
\[
 p(m)\sim\frac{\exp(c\sqrt m)}{4\sqrt3\,m},
 \qquad c=\pi\sqrt{2/3}.
\]
After adjusting an absolute constant to cover small values, for
$H/2\le j\le H$ this implies
\[
 \frac{p(j)}{p(H)}
 \le C\exp\!\left(-\frac{c(H-j)}{2\sqrt H}\right).
\]
The sum of $(2t+1)\exp(-ct/(2\sqrt H))$ over $t\ge0$ is $O(H)$.
The remaining part of \eqref{eq:Bdef} is at most
$O(H^2)p(\lfloor H/2\rfloor)$, which is $O(Hp(H))$ by the same
asymptotic. This proves the first inequality. The second follows
again from the partition asymptotic.
\end{proof}

\begin{lemma}[Euclidean logarithmic profile]
\label{lem:euclidean}
The finite quantity $W(a,b)$ itself, not only the minimum
two-block PT-rank, satisfies
\begin{equation}
\label{eq:W-upper}
 W(a,b)\le
 C(1+\log a+\log b)^3
 2^{\sqrt{(\log a)^2+(\log b)^2}}.
\end{equation}
The constant is independent of $a,b$ and of the coefficient field
or ring.
\end{lemma}

\begin{proof}
Set $\kappa=\pi\sqrt{2/3}/\log_e2$. The partition asymptotic first
gives $h(a)=O((1+\log a)^2)$. Minimality of $h(a)$, using
$p(h(a)-1)<a$ when $h(a)>1$, then gives
\begin{equation}
\label{eq:h-invert}
 \kappa\sqrt{h(a)}
 \le\log a+2\log(1+\log a)+C_0.
\end{equation}
The case $a=1$ is absorbed into the constant. Apply the same bound
to $h(b)$. If $x=\log a$ and $y=\log b$, the Euclidean triangle
inequality yields
\[
 \kappa\sqrt{h(a)+h(b)}
 \le\sqrt{x^2+y^2}
 +2\sqrt2\log(1+x+y)+O(1).
\]
Combine this with Lemma~\ref{lem:B-asymptotic}. The result has
power $2\sqrt2$ in the logarithmic prefactor; replacing it by three
gives \eqref{eq:W-upper}.
\end{proof}

\section{Orientation-preserving composition}

\subsection{Canonical orientations}

For $m\ge1$, let $\Ori_m$ consist of the subsets of $[m]$ that do
not contain the first coordinate. It has size $2^{m-1}$. Every
orientation is equivalent to one in $\Ori_m$ up to a full
transpose, which preserves rank.

\begin{definition}
A matrix on $m$ coordinate pairs has \emph{orientation bound $R$}
if it admits a decomposition
\begin{equation}
\label{eq:orientation-bound}
 M=\sum_{\alpha\in\Ori_m}A_\alpha^{\Gamma_\alpha},
 \qquad A_\alpha=U_\alpha V_\alpha^{\trsp},
\end{equation}
where every supplied factorization has inner dimension at most
$R$. Over a field this implies
$\PT(M)\le2^{m-1}R$.
\end{definition}

Keeping a matrix for each orientation is essential. Individual
rank-one summands from different orientations do not have
interchangeable row and column coordinates.

\subsection{Covariance under rectangular substitutions}

\begin{lemma}[Rectangular partial-transpose covariance]
\label{lem:covariance}
Let $U_A,V_A$ have the same number of rows and $a$ columns, and
let $U_B,V_B$ have the same number of rows and $b$ columns. For a
matrix $C$ on $[a]\times[b]$,
\begin{equation}
\label{eq:covariance}
 (U_A\otimes U_B)C^{\Gamma_B}(V_A\otimes V_B)^{\trsp}
 =\left[(U_A\otimes V_B)C(V_A\otimes U_B)^{\trsp}\right]^{\Gamma_B}.
\end{equation}
On the right, $\Gamma_B$ exchanges the entire physical $B$ block.
The identity is valid over every commutative unital ring.
\end{lemma}

\begin{proof}
The entry of either side at $(i_A,i_B;j_A,j_B)$ is
\[
 \sum_{p,q,r,s}
 U_A(i_A,p)U_B(i_B,q)
 C_{(p,s),(r,q)}
 V_A(j_A,r)V_B(j_B,s).
\]
Only commutative scalar arithmetic is used.
\end{proof}

\begin{lemma}[Generic orientation composition]
\label{lem:composition}
Suppose identities on $[a]\times[b]$ admit decompositions
$I_a\boxtimes I_b=C_0+C_1^{\Gamma_B}$ with supplied factorization
dimensions at most $W(a,b)$ for each $C_\epsilon$.
If matrices $A$ and $B$ on disjoint nonempty blocks of $k$ and
$\ell$ coordinate pairs have orientation bounds $R_A,R_B$, then
$A\boxtimes B$ has orientation bound at most $W(R_A,R_B)$.
A sum of at most $r$ such products has orientation bound at most
$rW(R_A,R_B)$.
\end{lemma}

\begin{proof}
Fix $\alpha\in\Ori_k$ and $\beta\in\Ori_\ell$. Pad the
factorizations of their matrices with zero columns so that
\[
 A_\alpha=U_AV_A^{\trsp},\qquad
 B_\beta=U_BV_B^{\trsp}
\]
have inner dimensions $R_A,R_B$. Insert
\[
 I_{R_A}\boxtimes I_{R_B}=C_0+C_1^{\Gamma_B}
\]
between these rectangular factors. The first term has a
factorization of dimension at most $W(R_A,R_B)$. By
Lemma~\ref{lem:covariance}, the second has such a factorization
after transposing the entire physical $B$ block.

After composing with the child orientations, the parent
orientations are
\begin{equation}
\label{eq:parent-orientations}
 \alpha\cup\beta
 \qquad\text{and}\qquad
 \alpha\cup\overline\beta,
\end{equation}
where the complement is taken inside the $B$ block. Every
canonical parent orientation identifies a unique child pair and
choice in \eqref{eq:parent-orientations}. Its restriction to $A$
is $\alpha$. Its first bit on $B$ says whether complementation
was used; undoing that complementation determines the canonical
$\beta$. Thus no two distinct child pairs contribute to the same
parent orientation within this one product.

For a sum of $r$ products, each parent orientation receives at
most one matrix from each product. Concatenating their
factorizations gives dimension at most $rW(R_A,R_B)$.
\end{proof}

\begin{remark}[What the uniqueness statement prevents]
The total number of orientations grows from
$2^{k-1}2^{\ell-1}$ to $2^{k+\ell-1}$. It does not multiply the
factorization dimension within an orientation. Losing this
distinction would replace the useful recurrence by one with an
unnecessary exponential factor at every level.
\end{remark}

\section{Paired coefficient tensors of arithmetic circuits}
\label{sec:circuits}

This section concerns a different polynomial encoding from the shifted
tensor used by Rossman and Zhu.  For an $n^d\times n^d$ matrix $M$, define
its \emph{paired coefficient polynomial}
\begin{equation}
\label{eq:paired-polynomial}
 F_M(Z)=\sum_{\mathbf i,\mathbf j\in[n]^d}
 M_{\mathbf i,\mathbf j}\prod_{t=1}^d Z_t(i_t,j_t),
\end{equation}
where the $d$ variable blocks $Z_t=\{Z_t(a,b):a,b\in[n]\}$
have size $n^2$.  All circuit bounds in this section apply to
\eqref{eq:paired-polynomial}, or to a specified squarefree
block component of a polynomial.  They do not identify
$F_M$ with $\operatorname{ShiftedTensor}(M)$.

Circuits are division-free commutative arithmetic circuits with binary
addition and multiplication gates; field or ring constants are allowed.
Their size counts all gates. In a homogeneous circuit we retain
syntactic degree labels: addition preserves the declared degree and
multiplication adds degrees. A zero polynomial may carry any declared
degree. The total-degree homogenization used below supplies these
labels, and restrictions preserve them. A syntactically set-multilinear circuit
assigns each gate a block scope $U\subseteq[d]$: constants have scope
$\varnothing$, a variable from $Z_t$ has scope $\{t\}$, addition
preserves scope, and a product combines disjoint scopes.  Known zero
components can be omitted; a zero output causes no difficulty.

\subsection{A balanced frontier with shared circuit contexts}

\begin{lemma}[Balanced circuit frontier]
\label{lem:circuit-frontier}
There is an absolute constant $C_0$ with the following property.
Let $f$ be the output of a homogeneous circuit of size $s$ and
degree $m\ge2$.  It has a decomposition
\begin{equation}
\label{eq:circuit-frontier}
 f=\sum_{u\in\mathcal F}g_u h_u,
 \qquad |\mathcal F|\le s,
\end{equation}
where both factors are homogeneous, their degrees sum to $m$,
each degree lies in $[m/3,2m/3]$, and each factor has a homogeneous
circuit of size at most $C_0s$.

If the given circuit is syntactically set-multilinear on a scope $U$,
the factor circuits may also be chosen syntactically set-multilinear.
For each summand their scopes form a partition of $U$, though this
partition may differ between summands.
\end{lemma}

\begin{proof}
Prune gates not used by the output.  Mark the gates of degree greater
than $2m/3$.  A marked addition gate has marked children.  A marked
multiplication gate has at most one marked child, because every gate
on a path to the output has degree at most $m$.

Starting with a context $q_{\rm out}=1$, process marked gates in
reverse topological order.  At a marked sum $v=v_0+v_1$, send
$q_v$ to each child.  At a marked product $v=v_0v_1$ with one
marked child, send $q_vf_{v_1}$ to $v_0$ if $v_0$ is marked,
or the symmetric expression if $v_1$ is marked.  Add all context
contributions arriving at one marked gate.  Stop at marked products
whose two children are unmarked; these gates form $\mathcal F$.
Nonzero input gates have degree at most one and are unmarked when
$m\ge2$.  An explicitly zero homogeneous input assigned a positive
formal degree contributes zero and may be discarded.

Successively substituting the marked-gate equations gives the exact
identity
\[
 f=\sum_{u\in\mathcal F}q_u f_{u_0}f_{u_1}.
\]
This is an identity in the coefficient ring.  It does not divide by
a path count, and shared paths are combined by addition in $q_u$.
Each $q_u$ is homogeneous of degree $m-\deg u$.

The circuit has at most $2s$ input wires.  Every propagated product
uses one new multiplication gate, and all context accumulations use
$O(s)$ new addition gates in total.  The original circuit together
with all contexts therefore has size $O(s)$.  In particular, a
context is not expanded into its possibly many paths.

For a frontier gate $u$, both children have degree at most $2m/3$,
while their degrees sum to more than $2m/3$.  Choose a child of
degree at least $m/3$ as $g_u$, and let $h_u$ be the other child
multiplied by $q_u$.  Both resulting degrees lie in $[m/3,2m/3]$,
and their circuits have size at most $C_0s$ for an absolute $C_0$.

In the set-multilinear case, induction on the reverse propagation
shows that every context arriving at a gate $v$ has scope
$U\setminus U_v$.  At a product, the sibling scope is disjoint
from this context scope; at an addition, incoming contexts have
the same scope.  Thus the context circuits and final factors
remain syntactically set-multilinear with the asserted complementary
scopes.
\end{proof}

\subsection{Varying partitions preserve the orientation bound}

The uniqueness in Lemma~\ref{lem:composition} is a statement for
each individual product.  It does not require different summands
to use the same partition of the physical pairs.  For one product
on disjoint scopes $A,B$, designate as $A$ the side containing the
first coordinate of the current scope.  A canonical parent
orientation restricts to a canonical orientation on $A$; its first
bit on $B$ determines whether the canonical orientation on $B$
was complemented.  Hence one product still contributes exactly
one coefficient matrix to each parent orientation.  A sum of $t$
such products has orientation width at most
\[
 t\max W(R_A,R_B),
\]
even when the partitions $A\mid B$ vary with the summand.

\begin{theorem}[Set-multilinear circuit bound]
\label{thm:setmultilinear-circuit}
Suppose $F_M$ has a syntactically set-multilinear arithmetic circuit
of size $s$.  There is an absolute constant $C$ such that, over
every field,
\begin{equation}
\label{eq:setmultilinear-circuit}
 \PT(M)\le2^{d-1}(2nsd)^{C\sqrt d}.
\end{equation}
Over every commutative unital ring the same bound holds for the number of
partial transposes of outer products in a supplied decomposition.
\end{theorem}

\begin{proof}
Apply Lemma~\ref{lem:circuit-frontier} recursively until each
factor has a singleton block scope.  The degree decreases by at
least a factor $2/3$ at every step, so the depth is $O(\log d)$.
The size of an individual descendant circuit increases by at most
$C_0$ per step.  Consequently every descendant has size at most
\[
 S_* = s d^{C_1}
\]
for an absolute $C_1$, after adjusting a constant for the bounded
small values of $d$.  This bounds the size of each circuit used
at a node; it does not claim that all recursively produced circuits
together have polynomial total size.

At a singleton scope, the coefficient matrix is an arbitrary
$n\times n$ matrix and has supplied factor width at most $n$.
At an internal node, each frontier term has two balanced scopes,
and the preceding varying-partition observation applies.  By
Lemma~\ref{lem:euclidean}, the logarithmic width satisfies a
numerical bound of the form
\begin{equation}
\label{eq:circuit-width-recurrence}
 L_v\le\sqrt{L_{v_0}^2+L_{v_1}^2}
       +\log S_*+C_2\log(2+L_{v_0}+L_{v_1}),
\end{equation}
where each child size is between one third and two thirds of the
parent size.  When several terms have different child sizes, choose
one attaining the largest numerical bound.  Repeating this choice
produces a balanced numerical tree with exactly $d$ singleton
leaves.  This tree bounds the maximum; the factor $S_*$ in the
recurrence retains all algebraic summands.

For completeness, put $\lambda=1+\log n+\log S_*$.  A coarse
induction gives $L_v\le K m_v\lambda$.  Indeed
$\sqrt{m_{v_0}^2+m_{v_1}^2}\le(\sqrt5/3)m_v$, so a sufficiently
large absolute $K$ absorbs the logarithmic errors.  Substituting
this coarse bound only into the final logarithm of
\eqref{eq:circuit-width-recurrence} bounds its additive error by
\[
 O\bigl(\log S_*+\log(m_v+1)+\log(2+\lambda)\bigr).
\]
In a band $2^j\le m_v<2^{j+1}$, balance gives $O(1)$ antichains,
each with at most $d/2^j$ nodes.  The Euclidean norm of one band's
error is therefore at most
\[
 O\!\left(\sqrt{d/2^j}
       [\log S_*+j+1+\log(2+\lambda)]\right).
\]
Summing the convergent geometric series, and adding the leaf norm
$\sqrt d\log n$, proves
$L_{\rm root}=O(\sqrt d\log(2nsd))$.  Sum the widths over the
$2^{d-1}$ canonical orientations.  All constructions use supplied
factorizations over the coefficient ring.
\end{proof}

\subsection{Lazy squarefree projection of ordinary circuits}

For a block set $U$, let $\mathsf{SF}_U(f)$ denote the component of
$f$ that has degree one in every block in $U$ and degree zero in
all other blocks.  If homogeneous polynomials $g,h$ have degrees
$k$ and $m-k$, respectively, then
\begin{equation}
\label{eq:lazy-squarefree}
 \mathsf{SF}_U(gh)
 =\sum_{\substack{A\subseteq U\\|A|=k}}
       \mathsf{SF}_A(g|_A)\,
       \mathsf{SF}_{U\setminus A}(h|_{U\setminus A}),
 \qquad |U|=m.
\end{equation}
Here $g|_A$ sets all variables outside the indicated blocks to zero.
The identity follows by assigning the $k$ block variables of a
contributing monomial to its factor $g$.  Homogeneity forces exactly
$k$ blocks on that side.  Terms with a repeated block, or a variable
outside $U$, cannot contribute.  This distribution of homogeneous
components is standard; see, for example, the component-distribution
argument in \cite[Lemma~9]{brandkaskiwang}.  The contribution here is
its balanced recursive use with canonical PT orientations and the
Euclidean crossing profile, which charges the total subset loss by
$O(d)$.

The essential point is that the child circuits in
\eqref{eq:lazy-squarefree} are just ordinary homogeneous circuits
with some inputs set to zero.  Their squarefree components are
tracked recursively; they are not first expanded into large
set-multilinear circuits.

\begin{theorem}[Ordinary circuit bound]
\label{thm:ordinary-circuit}
Let $f$ have a division-free arithmetic circuit of size $s$, and let
$M$ be its paired coefficient matrix, defined by
\[
 M_{\mathbf i,\mathbf j}
 =\left[\prod_{t=1}^d Z_t(i_t,j_t)\right]f.
\]
There is an absolute constant $C$ such that, over every field,
\begin{equation}
\label{eq:ordinary-circuit}
 \PT(M)\le 2^{Cd}(2nsd)^{C\sqrt d}.
\end{equation}
In particular this applies to any circuit computing $F_M$.
The corresponding supplied outer-product-length bound holds over
every commutative unital ring.
\end{theorem}

\begin{proof}
First retain homogeneous components of total degrees $0,\ldots,d$
at every gate.  At a product, component $k$ is the sum of the
products of components $a$ and $k-a$.  This constructs a
homogeneous circuit for the degree-$d$ component using
$O(s(d+1)^2)$ gates.  It is a formal coefficient calculation over
the ring; no interpolation or division is required.

For a current pair $(U,F)$ with $|U|=m$ and $F$ homogeneous of
degree $m$, apply Lemma~\ref{lem:circuit-frontier} to $F$ and
then apply \eqref{eq:lazy-squarefree} to every frontier product.
The resulting paired coefficient matrix is a sum of at most
$S_*2^m$ products on complementary scopes of sizes in
$[m/3,2m/3]$.  Each child is represented by an ordinary homogeneous
circuit, restricted to its scope, whose size has grown by at most
a fixed constant.  As before, all descendant sizes are bounded by
$S_*=s d^{O(1)}$.

The logarithmic orientation-width recurrence now becomes
\begin{equation}
\label{eq:ordinary-width-recurrence}
 L_v\le\sqrt{L_{v_0}^2+L_{v_1}^2}
       +m_v+\log S_*+C_3\log(2+L_{v_0}+L_{v_1}).
\end{equation}
The previous coarse linear induction still applies.  The only
new error is $m_v$.  Its Euclidean contribution from the band
$2^j\le m_v<2^{j+1}$ is at most
\[
 O\!\left(\sqrt{d/2^j}\,2^j\right)
 =O(\sqrt d\,2^{j/2}).
\]
Summing through $j\le\lceil\log d\rceil$ gives $O(d)$.
The other errors contribute $O(\sqrt d\log(2nsd))$, exactly
as in Theorem~\ref{thm:setmultilinear-circuit}.  Thus
\[
 L_{\rm root}=O\bigl(d+\sqrt d\log(2nsd)\bigr).
\]
The $2^{d-1}$ orientations add another $O(d)$ to the logarithm,
which proves \eqref{eq:ordinary-circuit}.
\end{proof}

\begin{remark}[A simpler, weaker cross-check]
Projecting every original gate onto every squarefree block scope
would instead give a syntactically set-multilinear circuit of size
$O(s3^d)$: a product component on $U$ sums over disjoint
$A,B$ with $A\cup B=U$, and the total number of such ordered
pairs is $3^d$.  Theorem~\ref{thm:setmultilinear-circuit} would
then give an additional $O(d\sqrt d)$ in the logarithm.  The
lazy projection in \eqref{eq:lazy-squarefree} improves this to
$O(d)$ by charging the subset choices at their own recursion size.
\end{remark}

\subsection{Border circuits and geometric border PT-rank}

\begin{proposition}[Border-circuit bound]
\label{prop:border-circuit}
Suppose $f$ lies in the coefficientwise Zariski closure, over an
algebraic closure of $\F$, of polynomials computed by division-free
circuits of size at most $s$.  Let $M$ be the matrix of its
coefficients with exactly one variable from every paired block.
Then
\[
 \overline{\PT}_{\F}(M)
 \le 2^{Cd}(2nsd)^{C\sqrt d}
\]
for an absolute constant $C$.
\end{proposition}

\begin{proof}
For fixed $s$, all circuit outputs have bounded total degree,
so they lie in a finite-dimensional coefficient space.  The map
that retains the specified paired coefficients is linear and
therefore Zariski-continuous.  By
Theorem~\ref{thm:ordinary-circuit}, its image on every size-$s$
circuit output belongs to the set of matrices with PT-rank at
most the displayed uniform bound.  Taking closures proves the
claim.  The theorem's raw-coefficient formulation is useful here:
the approximating outputs need not themselves be set-multilinear.
\end{proof}

\section{Uniform factors and their computational boundary}
The additional source input here is Lemma ``Explicit computation of
integral crossing coefficients'' in the canonical
\texttt{coefficient-computation.tex} of \cite{integralcrossing}.
For $h,E\ge1$, any entry of either supplied crossing factor is computable
in $2^{O(\sqrt{h+E})}$ bit operations and polynomial space in $h+E$;
the integer output and rational intermediates have polynomial bit length.
Generate the integer before specializing it to the target ring.
\subsection{Uniform factors from an ordinary circuit}

The coefficient-access lemma also makes the general circuit
bound constructive in the arithmetic model.  In this
subsection an effective coefficient ring means that its
addition and multiplication operations are supplied; their
cost is counted separately from integer constants generated
by the imported crossing algorithm.

\begin{theorem}[Constructive circuit factorization]
\label{thm:constructive-circuit}
Let $M$ be the paired squarefree-block coefficient matrix
of a division-free arithmetic circuit of size $s$, and put
\[
 H=d+\sqrt d\log(2nsd).
\]
Over every effective commutative unital ring, the decomposition
in Theorem~\ref{thm:ordinary-circuit} can be supplied with
at most $2^{O(H)}$ PT-basic terms such that each factor
entry is uniformly computable in $2^{O(H)}$ ring operations
and $2^{O(H)}$ bit operations for generating the integral
crossing constants.  A family of arithmetic circuits
representing all factor vectors can be generated with
total size and construction cost $2^{O(H)}$ in the same
model, using coordinate-indicator variables at the leaves.
\end{theorem}

\begin{proof}
Use the homogeneous circuit and balanced frontier of
Theorem~\ref{thm:ordinary-circuit}.  Every descendant
circuit has size at most $S_*=s d^{C_0}$.  Common deterministic
padding widths can be chosen by
\[
 R_1=n,\qquad
 R_m=S_*2^m
      \max_{m/3\le a\le2m/3}W(R_a,R_{m-a}),
\]
with integer rounding.  The earlier numerical-tree proof
gives
$\log R_m=O(m+\sqrt m\log(2nsd))$.
These widths and their partition-number parameters are
computable from the numerical recursion, independently of
circuit values.

A factor label records one frontier product, one subset
assigning blocks to its first factor, and one column in
the local crossing factor.  An orientation determines
the child orientations and whether the second child's
left and right factors are interchanged.  Invalid padded
labels denote zero columns.  These tests use circuit
syntax and integer labels, not polynomial identity testing.

For a fixed parent factor entry, its label fixes the
frontier term and the subset.  If its child widths are
$R_a,R_b$, the entry is a sum over $R_aR_b$ latent pairs
of products of two child entries and one explicit local
crossing coefficient.  The imported coefficient algorithm
computes that local coefficient in
$2^{O(\log R_a+\log R_b)}$ bit operations.  At a singleton,
factor the coefficient matrix as $A I$; an entry of $A$
is obtained by evaluating the degree-one child circuit at
the corresponding unit coordinate.

Writing $\lambda=\log(2nsd)$, a deliberately loose
worst-case cost therefore satisfies
\[
 T(m)\le 2^{O(m+\sqrt m\lambda)}
                   \bigl(T(a)+T(b)+1\bigr),
 \qquad a,b\in[m/3,2m/3].
\]
After taking logarithms, follow a larger-cost child.
Both $m$ and $\sqrt m$ decay geometrically on that path,
so $\log T(d)=O(H)$.  This estimate already charges all
latent-coordinate summations and requires no sharing
between descendant circuits on different block scopes.

Replacing a fixed physical coordinate at a leaf by its
indicator variables gives a circuit for the factor's
set-multilinear form.  Enumerating the $2^{d-1}$ orientations
and all $2^{O(H)}$ factor columns multiplies costs by at
most another $2^{O(H)}$, proving the total bound.
\end{proof}

The statement does not assume a polynomial bit overhead
for arbitrary integer input circuits.  Even a constant-only
circuit can repeatedly square $2$ and generate an integer
with exponentially many bits.  For a fixed finite field,
ordinary field-operation costs give a direct bit bound;
for other rings a bit-complexity assertion requires the
corresponding representation and value-size bounds.


\begin{thebibliography}{9}
\bibitem{integralcrossing}
icekylinx and GPT-6 Astra Pro.
\emph{Algorithm Frontier: Noncommutative Crossing}.
Repository snapshot
\texttt{2658b60c9d032611492792718bd9d39f721de09f},
October 28, 2026.
The relevant files are
\texttt{integral-pieces.tex},
\texttt{finite-crossing.tex}, and
\texttt{coefficient-computation.tex} under
\texttt{problems/noncommutative-crossing/manuscripts/}.
\url{https://github.com/icekylinx/algorithm-frontier/tree/2658b60c9d032611492792718bd9d39f721de09f/problems/noncommutative-crossing}.

\bibitem{rossmanzhu}
Benjamin Rossman and Davidson Zhu.
\emph{Multiquadratic Sum-of-Squares Lower Bounds Imply
$\mathrm{VNC}^1\ne\mathrm{VNP}$}.
ITCS 2026, LIPIcs 362, Article 113.
DOI: \href{https://doi.org/10.4230/LIPIcs.ITCS.2026.113}
{10.4230/LIPIcs.ITCS.2026.113}.
Full version:
\url{https://arxiv.org/html/2512.01227v1}.
The full proof of the ordered-ABP upper bound is in
Section 6.2.

\bibitem{anthropic}
Claude / Anthropic.
\emph{A $2^{o(n)}$ Time Algorithm for 3-SAT}.
October 26, 2026.
Repository snapshot
\texttt{16a5ce41304fcdaea4e0be7c132e1c361dc26da5},
file \texttt{eth/paper.pdf}; PDF blob
\texttt{77f4a8f043f65cecf0ee8e4024bf026aff1066e8}.
Appendix B supplies the underlying finite crossing-mode
construction credited by the integral extension.
\url{https://github.com/anthropics/formal-math/tree/16a5ce41304fcdaea4e0be7c132e1c361dc26da5/eth}.

\bibitem{dlmf}
NIST Digital Library of Mathematical Functions.
Section 27.14(iii), Equation 27.14.8: the
Hardy--Ramanujan asymptotic for the partition function.
\url{https://dlmf.nist.gov/27.14.E8}.


\bibitem{brandkaskiwang}
Cornelius Brand, Petteri Kaski, and Jiaheng Wang.
\emph{Partition Rank and Algebraic Circuit Lower Bounds}.
ESA 2026, LIPIcs 388, Article 63, pp.\ 63:1--63:13.
\url{https://doi.org/10.4230/LIPIcs.ESA.2026.63}.

\end{thebibliography}
\end{document}
